\documentclass[12pt,a4paper]{article}

\usepackage[margin=1in]{geometry}
\usepackage{amsmath,amssymb}
\usepackage{booktabs}
\usepackage{graphicx}
\usepackage{hyperref}
\usepackage{microtype}
\usepackage{xcolor}

\hypersetup{
  colorlinks=true,
  linkcolor=blue,
  citecolor=blue,
  urlcolor=blue
}

\title{Channel Estimation for Extremely Large-Scale Massive MIMO:\\
Far-Field, Near-Field, or Hybrid-Field?}
\author{Kelvin}
\date{\today}

\begin{document}

\maketitle

\begin{abstract}
Extremely large-scale massive multiple-input multiple-output (XL-MIMO)
systems provide high array gain and spatial resolution, but their large
apertures invalidate the conventional far-field plane-wave assumption for
many users. This report reviews channel estimation under far-field,
near-field, and hybrid-field propagation models. It summarizes the main
modeling assumptions, compares representative estimation approaches, and
discusses the trade-offs among accuracy, pilot overhead, and computational
complexity.
\end{abstract}

\tableofcontents
\newpage

\section{Introduction}

Massive MIMO improves spectral and energy efficiency by using a large
antenna array to serve multiple users simultaneously. In XL-MIMO systems,
the physical aperture becomes sufficiently large that some users or
scatterers may lie inside the array's Rayleigh distance. Their channels
therefore exhibit spherical wavefronts and depend on both angle and
distance, rather than angle alone.

This report focuses on the following questions:
\begin{itemize}
  \item When are far-field and near-field channel models appropriate?
  \item How can a system estimate channels containing both propagation types?
  \item What are the main performance and complexity trade-offs?
\end{itemize}

\section{System Model}

Consider a base station equipped with a uniform linear array of $M$
antennas. The received pilot signal can be written as
\begin{equation}
  \mathbf{y} = \mathbf{X}\mathbf{h} + \mathbf{n},
\end{equation}
where $\mathbf{X}$ is the pilot matrix, $\mathbf{h}$ is the channel vector,
and $\mathbf{n}$ is additive noise.

\subsection{Far-Field Model}

In the far field, the wavefront is approximated as planar. A path is
characterized primarily by its angle of arrival, and the steering vector
for a uniform linear array is
\begin{equation}
  \mathbf{a}_{\mathrm{FF}}(\theta)
  =
  \frac{1}{\sqrt{M}}
  \begin{bmatrix}
    1 &
    e^{-j 2\pi \frac{d}{\lambda}\sin\theta} &
    \cdots &
    e^{-j 2\pi \frac{d}{\lambda}(M-1)\sin\theta}
  \end{bmatrix}^{T}.
\end{equation}

\subsection{Near-Field Model}

In the near field, the spherical wavefront must be retained. The response
depends on both the angle $\theta$ and distance $r$:
\begin{equation}
  \left[\mathbf{a}_{\mathrm{NF}}(r,\theta)\right]_m
  =
  \frac{1}{\sqrt{M}}
  e^{-j\frac{2\pi}{\lambda}(r_m-r)},
\end{equation}
where $r_m$ is the distance from the source to antenna $m$.

\subsection{Hybrid-Field Model}

A practical XL-MIMO channel may contain far-field and near-field paths:
\begin{equation}
  \mathbf{h}
  =
  \sum_{\ell=1}^{L_{\mathrm{FF}}}
    \alpha_{\ell}\mathbf{a}_{\mathrm{FF}}(\theta_{\ell})
  +
  \sum_{k=1}^{L_{\mathrm{NF}}}
    \beta_{k}\mathbf{a}_{\mathrm{NF}}(r_k,\theta_k).
\end{equation}
This hybrid model motivates dictionaries or estimators that jointly
represent angular and polar-domain sparsity.

\section{Channel Estimation Methods}

\subsection{Least-Squares Estimation}

The least-squares estimate is
\begin{equation}
  \widehat{\mathbf{h}}_{\mathrm{LS}}
  =
  \left(\mathbf{X}^{H}\mathbf{X}\right)^{-1}
  \mathbf{X}^{H}\mathbf{y}.
\end{equation}
It is simple and does not require channel statistics, but it generally
requires substantial pilot overhead and can be sensitive to noise.

\subsection{Sparse Far-Field Estimation}

Far-field channels are often sparse in an angular discrete Fourier
transform dictionary. Compressed-sensing algorithms such as orthogonal
matching pursuit can exploit this sparsity to reduce pilot overhead.

\subsection{Sparse Near-Field Estimation}

Near-field estimation uses a polar-domain dictionary indexed by angle and
distance. This model improves accuracy for spherical wavefronts, at the
cost of a larger search space and higher dictionary coherence.

\subsection{Hybrid-Field Estimation}

Hybrid-field methods combine angular and polar-domain representations.
A practical estimator must determine whether each significant component
is better represented by a far-field or near-field atom while avoiding
duplicate or highly correlated selections.

\section{Comparison and Discussion}

\begin{table}[ht]
  \centering
  \caption{Qualitative comparison of propagation models.}
  \label{tab:comparison}
  \begin{tabular}{@{}llll@{}}
    \toprule
    Model & Parameters per path & Main advantage & Main limitation \\
    \midrule
    Far field & Angle & Low complexity & Near-field mismatch \\
    Near field & Angle and distance & Accurate spherical model & Large dictionary \\
    Hybrid field & Model, angle, distance & Broad applicability & Highest complexity \\
    \bottomrule
  \end{tabular}
\end{table}

The preferred model depends on array aperture, carrier wavelength, user
distance, scattering geometry, and available computational resources.
Model mismatch can produce a high error floor even when the signal-to-noise
ratio is large. Conversely, using a full near-field or hybrid-field model
for every user may introduce unnecessary complexity.

\section{Simulation Setup and Results}

Replace this section with the parameters and results used in the project.
A useful evaluation should report the normalized mean-square error
\begin{equation}
  \mathrm{NMSE}
  =
  \mathbb{E}\left[
    \frac{\lVert\mathbf{h}-\widehat{\mathbf{h}}\rVert_2^2}
         {\lVert\mathbf{h}\rVert_2^2}
  \right],
\end{equation}
as a function of signal-to-noise ratio, pilot length, user distance, or
the fraction of near-field paths.

% Example figure:
% \begin{figure}[ht]
%   \centering
%   \includegraphics[width=0.75\linewidth]{figures/nmse.pdf}
%   \caption{NMSE performance versus SNR.}
%   \label{fig:nmse}
% \end{figure}

\section{Conclusion}

XL-MIMO channel estimation requires propagation-aware modeling. Far-field
methods remain efficient for sufficiently distant users, while near-field
methods capture the distance-dependent structure introduced by spherical
wavefronts. Hybrid-field estimation provides a flexible solution when both
types of paths coexist, but its practical value depends on controlling
pilot overhead and computational complexity.

\begin{thebibliography}{9}

\bibitem{xlmimo_survey}
Insert the complete citation for the main XL-MIMO channel-estimation paper
used in this report.

\bibitem{massive_mimo}
T.~L. Marzetta,
``Noncooperative cellular wireless with unlimited numbers of base station
antennas,''
\emph{IEEE Transactions on Wireless Communications},
vol.~9, no.~11, pp.~3590--3600, 2010.

\end{thebibliography}

\end{document}
